\title{FlashAttention-3: Fast and Accurate Attention with Asynchrony and Low-precision}

\begin{abstract}
Attention mechanisms remain a fundamental component of the Transformer architecture, but also present a critical performance constraint in both large language models and applications requiring long-context processing. \textsc{FlashAttention} introduced efficient GPU-based attention computation by reducing memory operations. Nonetheless, it does not fully leverage the recent advancements available in the latest GPU hardware, as evidenced by \textsc{FlashAttention-2}, which only achieves 35\% hardware utilization on the H100 GPU. In this work, we introduce three innovations to accelerate attention on Hopper architecture GPUs: utilizing asynchronous Tensor Core and TMA operations to (1) concurrently perform computation and data transfer via warp specialization, (2) alternate between block-wise matrix multiplication and softmax calculation, and (3) employ block quantization alongside incoherent processing to make use of the hardware's FP8 low-precision capabilities. Our solution, \textsc{FlashAttention-3}, demonstrates 1.5--2.0$\times$ speed improvements on H100 GPUs, reaching 740 TFLOPs/s (75\% hardware utilization) with FP16, and approaching 1.2 PFLOPs/s using FP8. Experiments show that FP8 \textsc{FlashAttention-3} achieves 2.6$\times$ lower numerical error compared to a standard FP8 attention baseline.
\end{abstract}

\section{Introduction}
\label{sec:intro}

Within the Transformer architecture~\citep{vaswani2017attention}, the primary source of computational inefficiency arises from the attention mechanism, as evaluating self-attention between queries and keys incurs quadratic complexity relative to sequence length. Extending attention to handle longer contexts promises advancements such as the capacity to model and reason across multiple extended documents~\citep{guo2021longt5,shaham2022scrolls,peng2023yarn} or manage large-scale code repositories~\citep{roziere2023code, li2023starcoder}; it also facilitates the processing of additional modalities, including high-definition images~\citep{chen2022scaling}, audio~\citep{gulati2020conformer}, and video~\citep{ho2022video}, as well as novel applications like tracking lengthy user interaction histories~\citep{sun2019bert4rec} or supporting agents through extended workflows~\citep{yao2022react}. These prospects have motivated a broad effort toward expediting attention computation for lengthy sequences, employing techniques such as algorithmic approximation~\citep{katharopoulos2020transformers,choromanski2020rethinking,
tay2020efficient}, improvements in software implementations (\citep{rabe2021self,dao2022flashattention,kwon2023efficient}), and, in some instances, exploring alternative model architectures~\citep{peng2023rwkv,sun2023retentive,gu2023mamba}.

This study extends upon the contributions of \citet{dao2022flashattention}, who pioneered exact-attention algorithms informed by an in-depth understanding of GPU execution frameworks and underlying hardware. Dao et al. presented \textsc{FlashAttention}, an innovative approach utilizing a tiling method for attention computation that sidesteps multiple accesses to slow global memory by merging the entire suite of attention operations within a singular GPU kernel.

Building on this, \citet{dao2023flashattention2} proposed a revision of the algorithm, termed \textsc{FlashAttention-2}, which facilitates parallelism along the sequence length axis and organizes the forward pass's inner loop to process blocks of the key and value matrices, thus enhancing both GPU workload distribution and occupancy. Nevertheless, our findings reveal that \textsc{FlashAttention-2} underperforms in terms of resource utilization on recent GPU architectures when compared to well-tuned general matrix multiplication (GEMM) routines, demonstrating utilization rates of roughly 35\% versus 80-90\% on the Hopper H100 GPU. A contributing factor appears to be the reliance on instructions designed for the Ampere architecture, neglecting optimizations afforded by Hopper-specific instruction sets on Tensor Cores. Recent efforts, such as ThunkerKitten~\citep{spector2024thunder} and cuDNN 9~\citep{cudnn9}, illustrate that leveraging Hopper-tuned instructions and tile-centric implementation strategies can not only accelerate attention computation but also streamline code complexity.

At a more foundational level, the algorithm underlying \textsc{FlashAttention-2} operates within a basic synchronous paradigm and does not explicitly incorporate asynchrony or low-precision computation into its framework. Asynchrony, on the other hand, emerges from hardware optimizations that are engineered to expedite core operations in machine learning workloads—where specialized units focus on either matrix multiplication (like Tensor Cores) or memory access (e.g., Tensor Memory Accelerator -- TMA), distinct from general CUDA cores assigned to other logic, integer, or floating point tasks. The utilization of reduced-precision formats, such as FP8 on Hopper or FP4 on Blackwell—following the earlier adoption of FP16 (introduced with Pascal in 2017) and BF16 (in Ampere, 2020)—has reliably enabled a two- or even fourfold increase in throughput for a comparable allocation of power and silicon real estate. A discussion of the support provided by the Hopper architecture for these advancements is found in \cref{subsec:hardware}. A primary technical obstacle is to retrofit \textsc{FlashAttention-2} in order to leverage these hardware enhancements: achieving asynchrony entails simultaneous execution of matrix multiplication and softmax, even though a dependency exists between their outputs, while adopting low-precision arithmetic demands meticulous strategies to control quantization error, with particular attention to outlier feature values characteristic of LLMs~\citep{dettmers2208llm, sun2024massive}.

In response, we introduce \textsc{FlashAttention-3}, integrating and advancing three novel principles aimed at maximizing performance on the latest generation of GPUs:\footnote{The reported results focus on NVIDIA's Hopper architecture, yet our approach generalizes to any GPU platform with mature support for asynchronous execution and low-precision operations.}

\begin{enumerate}

\item \textbf{Producer-Consumer asynchrony:} We introduce a warp-specialized software pipelining approach that capitalizes on the decoupled execution capabilities of data transfer operations and Tensor Core computations. By allocating producers and consumers to different warps, this method increases the algorithm's potential to mask both memory access and instruction issuance delays.
\item \textbf{Hiding softmax under asynchronous block-wise GEMMs:} We achieve concurrency by overlapping softmax's low-throughput, non-GEMM computations—such as floating point multiply-adds and exponentials—with the asynchronous GEMM routines implemented via WGMMA instructions. To facilitate this, the \textsc{FlashAttention-2} algorithm is redesigned to remove certain sequential dependencies between the softmax computations and GEMM execution. For instance, in our two-stage scheme, the softmax step processes one block of the scores matrix while, in parallel, WGMMA carries out asynchronous computations on the subsequent block.
\item \textbf{Hardware-accelerated low-precision GEMM:} The forward pass procedure is revised to support FP8 Tensor Cores during GEMM operations, resulting in an observed near doubling of TFLOPs/s performance. Accommodating this enhancement necessitates reconciling the layout constraints imposed by WGMMA on FP32 accumulators and FP8 operand blocks in memory. To address accuracy degradation due to reduced FP8 precision, we apply block quantization and incoherent computation techniques.
\end{enumerate}

To empirically assess the effectiveness of our approach, we conduct benchmarks of \textsc{FlashAttention-3} on the H100 SXM5 GPU across various configurations. Our results demonstrate that (1) for the forward pass, the FP16 version delivers a speedup of 1.5 to 2.0$\times$ compared to \textsc{FlashAttention-2} (with performance peaking at 740 TFLOPs/s), and achieves a 1.5–1.75$\times$ speedup during the backward pass; (2) the FP8 variant attains nearly 1.2 PFLOPs/s; and (3) at longer sequence lengths, FP16 displays superior performance, while FP8 achieves competitive throughput\footnote{Specifically, with head dimension 64, FP8 in \textsc{FlashAttention-3} surpasses alternatives, while for head dimensions 128 and 256, it matches their performance in non-causal settings and falls slightly behind when causal masking is employed.} compared to a leading attention implementation from NVIDIA's cuDNN. Additionally, we verify that FP16 \textsc{FlashAttention-3} maintains the same level of numerical error as \textsc{FlashAttention-2}—and surpasses conventional attention mechanisms—since intermediate results (such as softmax rescaling) are handled in FP32 precision. Further, in scenarios with outlier features, FP8 \textsc{FlashAttention-3} using block quantization and incoherent processing yields 2.6$\times$ greater accuracy than traditional attention using per-tensor quantization.

\textsc{FlashAttention-3} is released under a permissive open-source license\footnote{\textsc{FlashAttention-3} is accessible at \texttt{https://github.com/Dao-AILab/flash-attention}}. We also intend to integrate the library with PyTorch and Hugging Face ecosystems to maximize its availability to the research and development communities.

\section{Background: Multi-Head Attention and GPU Characteristics}
\label{sec:background}

\subsection{Multi-Head Attention}
\label{subsec:multi_head_attn}

Consider the query, key, and value matrices $\mathbf{Q}, \mathbf{K}, \mathbf{V} \in \mathbb{R}^{N \times d}$ for a single attention head, where $N$ denotes the sequence length and $d$ specifies the dimension of each head. The attention mechanism generates the output $\mathbf{O}$ via the following steps:

\begin{equation*}
  \mathbf{S} = \alpha \mathbf{Q} \mathbf{K}^\top \in \mathbb{R}^{N \times N}, \quad \mathbf{P} = \mathrm{softmax}(\mathbf{S}) \in \mathbb{R}^{N \times N}, \quad \mathbf{O} = \mathbf{P}\mathbf{V} \in \mathbb{R}^{N \times d},
\end{equation*}

Here, $\mathrm{softmax}$ operates along each row, and it is standard to choose the scaling parameter as $\alpha = 1/\sqrt{d}$. To avoid numerical issues stemming from large exponentials, we subtract $\mathrm{rowmax}(\mathbf{S})$ from $\mathbf{S}$. In multi-head attention (MHA), each attention head utilizes individual projection matrices for queries, keys, and values; as a result, these calculations are performed concurrently across all heads and batches, ultimately yielding the complete output tensor.

Consider $\phi$ as a scalar loss function, and let us denote its gradient with respect to an argument by $\mathbf{d}(-) = \partial \phi / \partial (-)$. Starting from the gradient at the output, $\mathbf{dO} \in \mathbb{R}^{N \times d}$, the gradients with respect to $\mathbf{Q}$, $\mathbf{K}$, and $\mathbf{V}$ can be determined via the chain rule as shown below:

\begin{align*}
  \mathbf{dV} &= \mathbf{P}^\top \mathbf{dO} \in \mathbb{R}^{N \times d} \\
  \mathbf{dP} &= \mathbf{dO} \mathbf{V}^\top \in \mathbb{R}^{N \times N} \\
  \mathbf{dS} &= \mathrm{dsoftmax} (\mathbf{dP}) \in \mathbb{R}^{N \times N} \\
  \mathbf{dQ} &= \alpha \mathbf{dS} \mathbf{K} \in \mathbb{R}^{N \times d} \\
  \mathbf{dK} &= \alpha \mathbf{dS}^\top \mathbf{Q} \in \mathbb{R}^{N \times d},
\end{align*}

Specifically, for the softmax Jacobian applied to vector-valued $s$, we use the identity $\mathbf{d}s = (\mathrm{diag}(p) - p p^\top)\mathbf{d}p$, with $p = \mathrm{softmax}(s)$. The operator $\mathrm{dsoftmax}(\mathbf{dP})$ indicates that this transformation is applied to each row individually. This procedure is inherently parallelizable over both the different heads and individual elements within each batch for the backward phase of MHA.

\subsection{GPU hardware characteristics and execution model}
\label{subsec:hardware}

We describe the aspects of the GPU's execution model relevant for \textsc{FlashAttention-3}, with a focus on the NVIDIA Hopper architecture as a concrete instantiation of this model.

\paragraph{Memory hierarchy:} The GPU's memory system is structured in a layered hierarchy of data storage locations, where greater bandwidth is accompanied by reduced storage capacity (\cref{tab:gpu-hierarchy})\footnote{\citet{luo2024benchmarking} specifies a shared memory bandwidth of 128 bytes per clock cycle per SM. For our calculations, we multiplied this value by 132 SMs and the boost clock frequency of 1830 MHz.}. The primary large-scale storage, global memory (GMEM) or HBM, resides off-chip and is accessible to every streaming multiprocessor (SM). When data is accessed from GMEM, it is automatically cached in the on-chip L2 cache. Each SM is also equipped with its own small-sized, explicitly managed and highly banked cache referred to as shared memory (SMEM). At the top of the hierarchy is the register file, located within each SM.

\paragraph{Thread hierarchy:} GPUs structure their computation model by grouping execution entities known as threads. This hierarchy is organized from the smallest units upward: individual threads, then warps (each containing 32 threads), warpgroups (made up of four sequential warps), threadblocks (also known as cooperative thread arrays or CTAs), threadblock clusters (introduced with Hopper architecture), and, at the highest level, grids.

There is a strong correlation between these two hierarchies. Threads belonging to a particular CTA are jointly scheduled on a single SM, while multiple CTAs that belong to the same cluster are assigned to the same GPC. All threads inside a CTA can directly access SMEM, while each individual thread possesses up to 256 dedicated registers (RMEM) that are not shared with others.

\paragraph{Asynchrony and warp-specialization:}

GPUs function as throughput-oriented processors, leveraging parallelism and asynchronous operations to mask the delays associated with memory accesses and instruction execution. To facilitate asynchronous data transfers between GMEM and SMEM, Hopper introduces the Tensor Memory Accelerator (TMA), a specialized hardware component \cite[\S7.29]{cuda}. Additionally, differing from earlier architectures like Ampere, Hopper’s Tensor Core—accessible through the warpgroup-wide WGMMA instruction \cite[\S9.7.14]{ptx}—operates asynchronously and has the capability to directly fetch its input operands from shared memory.

The presence of hardware mechanisms enabling asynchrony facilitates the implementation of warp-specialized kernels, wherein the warps within a CTA are allocated to specific roles as either producers or consumers, thus restricting their activity to solely data transfer or computation. This specialization generally enhances the compiler's capability to produce highly efficient instruction schedules \citep{warp-specialization-2011}. Moreover, the Hopper architecture introduces functionality for reallocating register resources dynamically among warpgroups through the \verb|setmaxnreg| interface \cite[\S9.7.17.1]{ptx}, permitting warps engaged in MMA operations to access a greater portion of RMEM compared to those warps performing only TMA, which typically require participation from only one thread.

\paragraph{Low-precision number formats:}
\label{sec:low-precision-gpu}
Contemporary GPUs are equipped with dedicated hardware components to enhance the performance of low-precision arithmetic. For instance, the WGMMA instruction leverages FP8 Tensor Cores on Hopper GPUs, enabling a twofold increase in throughput per SM relative to computations performed in FP16 or BF16.

Nevertheless, to properly utilize FP8 WGMMA, one must be aware of the specific memory layout requirements imposed on its operands. Consider a GEMM operation that computes $A \times B^{\top}$, where $A$ is an $M \times K$ matrix and $B$ an $N \times K$ matrix. An operand (either $A$ or $B$) is said to be \emph{mn-major} if its data is stored contiguously along the outer $M$ or $N$ axis, and \emph{k-major} if the contiguous dimension is the inner $K$ axis instead. For FP16 WGMMA, the shared memory operands can be supplied in either mn-major or k-major arrangements. In contrast, FP8 WGMMA restricts inputs in shared memory exclusively to the k-major configuration. Additionally, certain applications, such as those requiring fused execution of sequential GEMMs (as in attention mechanisms), encounter complications: conflicting layouts between the FP32 accumulator and subsequent FP8 operands hinder seamless invocation of dependent FP8 WGMMAs within a unified kernel.

In the context of attention, these layout restrictions entail certain modifications to the design of an FP8 algorithm, which we describe in \cref{sec:algofp8}.

\subsection{Standard Attention and Flash Attention} In line with \citet{dao2022flashattention}, we refer to \textbf{standard attention} as the GPU-based attention computation that explicitly stores the intermediate matrices $\mathbf{S}$ and $\mathbf{P}$ in HBM. The core concept behind \textsc{FlashAttention} involves utilizing a local variant of the softmax reduction to eliminate the need for costly intermediate memory accesses, thereby merging the attention operation into a unified kernel. The localized softmax is executed in lines \ref{code-ws:softmax_start}-\ref{code-ws:softmax_end} of the consumer mainloop in \cref{alg:flash3_wgmma_ws_only}, along with the appropriate rescaling of sections of $\mathbf{O}$. A straightforward demonstration showing that this method accurately produces $\mathbf{O}$ can be found in \cite[\S 2.3.1]{dao2023flashattention2}.

\section{FlashAttention-3: Algorithm}
\label{sec:algo}

This section provides an overview of the \textsc{FlashAttention-3} algorithm. To streamline the discussion, our attention is directed to the forward computation, while the procedure for the backward pass can be found in~\cref{sec:algo_ws_bwd}. We begin by outlining the incorporation of warp-specialization, coupled with a circular SMEM buffer, into the foundational framework established by \textsc{FlashAttention-2}. Next, we elaborate on the strategy of leveraging WGMMA asynchrony to construct a two-stage pipeline that overlaps GEMM with the softmax operation. Lastly, we detail the necessary adjustments for FP8, covering aspects of layout compatibility as well as accuracy enhancement through block quantization and processing with incoherence.

\subsection{Producer-Consumer asynchrony through warp-specialization and pingpong scheduling}
\label{sec:algo_ws}

\paragraph{Warp-specialization}
Similar to the design of \textsc{FlashAttention-2}, the forward computation in \textsc{FlashAttention-3} is highly parallelizable across the dimensions of batch size, number of attention heads, and the length of the query sequence. Consequently, it is sufficient to present the algorithm at the CTA level, where each CTA processes a specific tile $\mathbf{Q}_i$ of the query matrix and produces the matching output tile $\mathbf{O}_i$. For clarity, we first explain the warp-specialization approach using a circular SMEM buffer, intentionally excluding any overlap between GEMM and softmax computations. Let $d$ represent the head dimension and $N$ denote the sequence length. By selecting a block size $B_r$ for the queries, $\mathbf{Q}$ can be partitioned into $T_r = \lceil \frac{N}{B_r} \rceil$ blocks labeled as $\mathbf{Q}_1, \ldots, \mathbf{Q}_{T_r}$.

\begin{algorithm}[H]
    \caption{\small\label{alg:flash3_wgmma_ws_only}\textsc{FlashAttention-3} forward pass \textbf{excluding} intra-consumer overlap -- CTA perspective}
    \begin{algorithmic}[1]
\REQUIRE Matrices $\mathbf{Q}_i \in \mathbb{R}^{B_r \times d}$ and $\mathbf{K}, \mathbf{V} \in \mathbb{R}^{N \times d}$ stored in HBM; key block dimension $B_c$ and $T_c = \lceil \frac{N}{B_c} \rceil$.
\STATE Set up a pipeline structure to oversee barrier synchronization, leveraging an $s$-stage circular SMEM buffer.
\IF {operating in producer warpgroup}
\STATE Release a specified quantity of registers.
\STATE Trigger the load of $\mathbf{Q}_i$ from HBM into shared memory.
\STATE Once loading completes, notify the consumers regarding the availability of $\mathbf{Q}_i$.
\FOR{$0 \le j < T_c$}
    \STATE Await consumption of the $(j\,\%\,s)$th stage of the buffer.
    \STATE Initiate loads of $\mathbf{K}_j, \mathbf{V}_j$ from HBM into shared memory at the buffer’s $(j\,\%\,s)$th slot.
    \STATE Upon finishing, signal to consumers that $\mathbf{K}_j, \mathbf{V}_j$ are now ready.
\ENDFOR
\ELSE \STATE Allocate a specified number of registers as determined by the count of consumer warps.
\STATE Locally, set $\mathbf{O}_i = (0) \in \mathbb{R}^{B_r \times d}$ along with $\ell_i, m_i = (0), (-\infty) \in \mathbb{R}^{B_r}$.
\STATE Wait until $\mathbf{Q}_i$ has been transferred into shared memory.
\FOR{$0 \le j < T_c$}
\STATE Hold for $\mathbf{K}_j$ to become available in shared memory.
\STATE Compute $\mathbf{S}_i^{(j)} = \mathbf{Q}_i \mathbf{K}_j^T$ using SS-GEMM. Commit and block. \label{alg:ws_only_gemm1}
\STATE Save $m_i^{\mathrm{old}} = m_i$, then update $m_i = \mathrm{max}(m_i^{\mathrm{old}}, \mathrm{rowmax}(\mathbf{S}_i^{(j)}))$. \label{code-ws:softmax_start}
\STATE Evaluate $\widetilde{\mathbf{P}}_i^{(j)} = \mathrm{exp}(\mathbf{S}_i^{(j)} - m_i)$ and $\ell_i = \mathrm{exp}(m_i^{\mathrm{old}} - m_i)\ell_i + \mathrm{rowsum}(\widetilde{\mathbf{P}}_i^{(j)})$. \label{code-ws:softmax_end}
\STATE Wait for the arrival of $\mathbf{V}_j$ in shared memory.
\STATE Update $\mathbf{O}_i = \mathrm{diag}(\mathrm{exp}(m_i^{\mathrm{old}} - m_i))^{-1} \mathbf{O}_i + \widetilde{\mathbf{P}}_i^{(j)} \mathbf{V}_j$ via RS-GEMM. Commit and block. \label{alg:ws_only_gemm2}
\STATE Release the buffer’s $(j\,\%\,s)$th stage to allow producer progression.
\ENDFOR
\STATE Compute $\mathbf{O}_i = \mathrm{diag}(\ell_i)^{-1}\mathbf{O}_i$ as well as $L_i = m_i + \log(\ell_i)$.
\STATE Output $\mathbf{O}_i$ and $L_i$ to HBM as the $i$th portion of $\mathbf{O}$ and $L$.
\ENDIF
\end{algorithmic}
\end{algorithm}

In our implementation of \cref{alg:flash3_wgmma_ws_only} on Hopper, we employ \verb|setmaxnreg| to handle (de)allocation tasks, utilize TMA to load both $\mathbf{Q}_i$ and $\{ \mathbf{K}_j, \mathbf{V}_j \}_{0 \leq j < T_c}$, and make use of WGMMA for performing the GEMM computations in the consumer mainloop. The prefixes SS or RS denote whether the initial operand is being sourced from shared memory or from the register file, respectively. When examining the execution pattern described in \cref{alg:flash3_wgmma_ws_only}, it is important to observe that TMA load instructions execute asynchronously, meaning their initiation does not block on the completion of prior loads. Additionally, during the producer mainloop, for the initial $s$ iterations—corresponding to the phase where the buffer is being initially populated—no waits are required.

\paragraph{Pingpong scheduling} Leveraging the asynchronous execution of WGMMA and TMA, together with warp specialization, provides an opportunity to concurrently process the softmax for one warpgroup while another warpgroup is performing a GEMM operation. To understand the rationale, it is important to note the significant disparity in throughput between matmul and non-matmul tasks on contemporary hardware accelerators. For instance, the H100 SXM5 GPU is capable of 989 TFLOPS in FP16 matmul operations, yet delivers only 3.9 TFLOPS for special function calculations such as exponentials\footnote{According to the CUDA programming guide, each streaming multiprocessor (SM) can perform 16 special function operations per clock cycle. By multiplying 16 operations by 132 SMs and a 1830 MHz clock rate, we derive the 3.9 TFLOPS figure for special functions.}, which are essential for softmax. In an FP16 attention forward pass with a head dimension of 128, there are 512 times as many matmul FLOPS as exponential operations, but the hardware executes exponentials at a rate 256 times lower, meaning the exponential operations could account for as much as 50\% of the total computation time relative to matmul. The imbalance grows more pronounced in FP8, where matmul throughput doubles while exponential performance remains unchanged.

Given that the exponential operation is handled by a dedicated hardware component—the multi-function unit—it is optimal to overlap the exponential computation with the period during which the Tensor Cores are executing the matrix multiplication. To achieve this concurrency, we implement synchronization barriers (using \texttt{bar.sync} instructions) to ensure the GEMMs (GEMM1: $\mathbf{P} \mathbf{V}$ from the current iteration, and GEMM0: $\mathbf{Q} \mathbf{K}^\top$ from the subsequent iteration) executed by warpgroup 1 occur prior to the GEMMs in warpgroup 2. This mechanism enables the softmax operation for warpgroup 1 to be performed in parallel with the GEMMs of warpgroup 2. The execution order is then alternated, so that warpgroup 2 processes softmax concurrently as warpgroup 1 initiates its own GEMMs, following a ``pingpong'' schedule. An illustration of this scheduling technique is presented in~\cref{fig:pingpong_scheduling}. Although the overlap achieved via pingpong scheduling may not be perfectly aligned as depicted, we consistently observe notable performance enhancements in practice (such as an increase from 570 TFLOPS to 620–640 TFLOPS on FP16 forward passes with head dimension 128 and a sequence length of 8192).

\paragraph{Attention variants} For multi-query attention~\citep{shazeer2019fast} and grouped query attention~\citep{ainslie2023gqa}, we adopt the tensor indexing modifications introduced in \textsc{FlashAttention-2} to circumvent unnecessary replication of $\mathbf{K}$ and $\mathbf{V}$ within HBM.

\subsection{Intra-warpgroup overlapping GEMMs and softmax}

Even within one warpgroup, we can overlap some instructions in the softmax with some instructions in the GEMMs. We describe one technique to do so.

Within the attention mechanism, operations inside the main (inner) loop exhibit inherent sequential dependencies, which limit opportunities for parallel execution in each individual iteration. For instance, the computation of the (local) softmax (refer to lines \ref{code-ws:softmax_start} through \ref{code-ws:softmax_end}) is contingent on the outcome $\mathbf{S}_i^{(j)}$ produced by the initial GEMM, while the subsequent GEMM then requires $\widetilde{\mathbf{P}}_i^{(j)}$ as its input. The presence of explicit wait instructions at lines \ref{alg:ws_only_gemm1} and \ref{alg:ws_only_gemm2} in \cref{alg:flash3_wgmma_ws_only} further enforces a serialized execution pattern between the softmax and GEMM operations. Nevertheless, it is feasible to mitigate these serial dependencies by introducing pipelining across loop iterations, achieved via supplementary buffer allocation in registers. Based on this principle, we introduce a two-stage\footnote{It is important to observe that the count of stages in the overlapping scheme is limited by the number $s$ of stages in the circular SMEM buffer, but they need not coincide.} pipelined GEMM-softmax algorithm:

\begin{algorithm}[H]
    \caption{\small\label{alg:flash3_wgmma}\textsc{FlashAttention-3} Consumer Warpgroup: Forward Pass Procedure}
    \begin{algorithmic}[1]
      \REQUIRE  Input: matrices $\mathbf{Q}_i \in \mathbb{R}^{B_r \times d}$ and $\mathbf{K}, \mathbf{V} \in \mathbb{R}^{N \times d}$ stored in HBM, key block dimension $B_c$, with total number of blocks $T_c = \lceil \frac{N}{B_c} \rceil$.
      \STATE Reassign registers in advance depending on the count of consumer warps.
      \STATE Allocate and initialize on-chip: set $\mathbf{O}_i = (0) \in \mathbb{R}^{B_r \times d}$, and both $\ell_i = (0)$, $m_i = (-\infty) \in \mathbb{R}^{B_r}$.
      \STATE Ensure $\mathbf{Q}_i$ and initial $\mathbf{K}_0$ are present in shared memory.
      \STATE Use WGMMA to evaluate $\mathbf{S}_{\mathrm{cur}} = \mathbf{Q}_i \mathbf{K}_0^T$. Commit the operation and synchronize.
      \STATE Free up stage $0$ in the $\mathbf{K}$ buffer.
      \STATE Derive values for $m_{i}$, $\tilde{\mathbf{P}}_{\mathrm{cur}}$, and $\ell_{i}$ from $\mathbf{S}_{\mathrm{cur}}$, then update $\mathbf{O}_i$ accordingly.
      \FOR{$1 \le j < T_c - 1$}
        \STATE Await availability of $\mathbf{K}_j$ in shared memory.
        \label{code:mainloop_start}
        \STATE Compute $\mathbf{S}_{\mathrm{next}} = \mathbf{Q}_i \mathbf{K}_{j}^T$ by means of WGMMA; submit for execution without blocking. \label{code:first_wgmma}
        \STATE Wait for shared memory to contain $\mathbf{V}_{j-1}$.
        \STATE With WGMMA, accumulate $\mathbf{O}_{i} = \mathbf{O}_{i} + \tilde{\mathbf{P}}_{\mathrm{cur}} \mathbf{V}_{j-1}$; commit without synchronization. \label{code:second_wgmma}
        \STATE Ensure completion of the WGMMA involving $\mathbf{Q}_i \mathbf{K}_{j}^T$.
        \STATE Based on $\mathbf{S}_{\mathrm{next}}$, compute $m_{i}$, $\tilde{\mathbf{P}}_{\mathrm{next}}$, and $\ell_{i}$. \label{code:softmax}
        \STATE Upon completion of the WGMMA involving $\tilde{\mathbf{P}}_{\mathrm{cur}} \mathbf{V}_{j-1}$, normalize $\mathbf{O}_i$.
        \STATE Free stages $(j\,\%\,s)$ and $(j-1\,\%\,s)$ from $\mathbf{K}$ and $\mathbf{V}$ buffers, respectively.
        \STATE Copy contents of $\mathbf{S}_{\mathrm{next}}$ to $\mathbf{S}_{\mathrm{cur}}$.
        \label{code:mainloop_end}
      \ENDFOR \STATE Wait for transfer of $\mathbf{V}_{T_c - 1}$ to shared memory to complete.
      \STATE Calculate
      $\mathbf{O}_{i} = \mathbf{O}_{i} + \tilde{\mathbf{P}}_{\mathrm{last}} \mathbf{V}_{T_c - 1}$ leveraging WGMMA, committing and waiting for completion.

\STATE Epilogue: Adjust $\mathbf{O}_{i}$ by applying the scaling factor $m_{i}$. Determine $L_{i}$ using both $m_{i}$ and $\ell_{i}$.  
      Store the results of $\mathbf{O}_{i}$ and $L_{i}$ into HBM, assigning them as the $i$-th segments of $\mathbf{O}$ and $L$ respectively.  
    \end{algorithmic}
  \end{algorithm}

  \cref{alg:flash3_wgmma} substitutes for the consumer path presented in \cref{alg:flash3_wgmma_ws_only}, thereby constructing the full \textsc{FlashAttention-3} algorithm when operating with FP16 data. From an abstract perspective, we treat WGMMA as an equivalent for asynchronous GEMM. During the principal loop (spanning lines \ref{code:mainloop_start} through \ref{code:mainloop_end}), the execution of the second WGMMA instruction for iteration $j$ (refer to line \ref{code:second_wgmma}) proceeds concurrently with the softmax computations pertaining to iteration $j+1$ (see line \ref{code:softmax}).

Although the aforementioned pipelined architecture has the potential for significant theoretical speedup, its actual implementation introduces several practical challenges:
\paragraph{Compiler reordering} The pseudocode assumes an ideal, sequential order of execution; however, the compiler (NVCC) may restructure instructions to optimize performance. Such reordering can compromise the intentionally designed interleaving of WGMMA and other operations, possibly resulting in performance degradation or unanticipated side effects. Examination of the compiled SASS code indicates that NVCC generally preserves the intended instruction overlap (see Section~\ref{sec:2-stage-sass}).
\paragraph{Register pressure} Achieving peak efficiency requires limiting register spilling as much as possible. Nonetheless, the 2-stage pipeline inherently demands additional registers to hold intermediate values and retain the context of each stage. In particular, registers must accommodate an extra $\mathbf{S}_{\mathrm{next}}$, incurring an overhead of $B_r \times B_c \times \text{sizeof}(\text{float})$ extra bytes per threadblock. This greater register requirement may directly compete with the pursuit of larger block sizes—a strategy that itself increases register usage. Thus, balancing these aspects often necessitates empirical tuning guided by profiling outcomes.
\paragraph{3-stage pipelining} Building upon the 2-stage scheme, we introduce a 3-stage version that overlaps the second WGMMA phase with the softmax computation. This extension holds promise for better Tensor Core utilization, yet comes at the cost of even more registers due to the extra pipeline stage. As a result, the tension between maximizing tile size and increasing pipeline depth becomes more pronounced. A comprehensive explanation and evaluation of the 3-stage strategy are available in ~\cref{sec:3-stage}.

\subsection{Low-precision with FP8}
\label{sec:algofp8}

\textbf{Efficiency: layout transformations.} Executing the forward computation in \textsc{FlashAttention-3} using FP8 precision introduces layout conformance difficulties absent in the FP16 setting.

Generally, input tensors $\mathbf{Q}$, $\mathbf{K}$, and $\mathbf{V}$ are organized in memory such that the head dimension is contiguous. However, for the second GEMM involving FP8 WGMMA, the k-major requirement mandates that $\mathbf{V}$, or more precisely, its tiles loaded into SMEM, be stored contiguously along the sequence length dimension. Since a TMA load operation cannot alter the axis of contiguity, one must either (1) execute a pre-processing transpose of $\mathbf{V}$ in GMEM, or (2) perform an in-kernel transpose on tiles of $\mathbf{V}$ after they are resident in SMEM. Option (1) itself can be realized by either (1a) integrating the transpose into the epilogue of a prior computation step, for instance during rotary embedding, or (1b) utilizing a dedicated pre-processing kernel that transposes the data, switching the stride order between the sequence length and head dimensions.\footnote{A specialized transpose kernel can attain bandwidth close to the hardware’s theoretical limit~\citep{colfax_cutlass_transpose_2024}.} Despite this, strategy (1a) is not easily adaptable for inclusion in general-purpose libraries, and approach (1b) results in excessive memory usage, which is suboptimal especially during inference when memory bandwidth is often the limiting factor.

For FP8 \textsc{FlashAttention-3}, we select approach (2). To carry out the in-kernel transpose, we exploit the LDSM (\verb|ldmatrix|) and STSM (\verb|stmatrix|) instructions, whereby a warp of threads collaboratively loads data from SMEM to RMEM and stores from RMEM to SMEM with a transfer granularity of 128 bytes.\footnote{The PTX specification states that LDSM and STSM perform $8 \times 8$ matrix operations using 16-bit elements \cite[\S 9.7.13.4.15-16]{ptx}. Nevertheless, by packing pairs of 8-bit values, it is feasible to utilize LDSM/STSM for FP8 data. Notably, the transpose operations for LDSM/STSM cannot decompose packed 8-bit items, leading to the need for some register moves between the LDSM and STSM steps in order to enable tile-level transposition; the specifics are omitted here.} LDSM/STSM are both efficient with respect to register usage, making it possible to run them directly in the producer warpgroup, and are able to transpose data layouts during the memory movement. Furthermore, following the initial iteration, it is possible to schedule the transpose of the subsequent $\mathbf{V}$ tile so that it overlaps with the two WGMMAs corresponding to the previous $\mathbf{V}$ and the current $\mathbf{K}$ tile.

Second, we note that, in contrast to FP16, the arrangement in memory of the FP32 accumulator resulting from an FP8 WGMMA operation does not match the memory layout of operand A when stored in registers. This distinction is illustrated by partial examples in \cref{fig:rmem_accum} and \cref{fig:rmem_operand}, where the sequence denotes the order in which threads retain the data in registers. By employing byte permutation instructions, the output of one WGMMA's accumulator can be reshaped to match the required format for a subsequent WGMMA, as well as align with the structure of the $\mathbf{V}$ tile after an in-kernel transpose. To be specific, according to \cref{fig:rmem_accum}, we reorder the bytes to the following sequence:
$$\{ \verb|d0 d1 d4 d5 d2 d3 d6 d7| \},$$
and apply this permutation consistently for each consecutive 8-byte segment. This operation effectively permutes the columns of the logical $\mathbf{P}$ tile (for instance, columns $0189$ are brought to the foremost positions). For the subsequent WGMMA to generate the correct output tile, the in-kernel transpose can be designed to output the $\mathbf{V}$ tile with a corresponding row permutation.\footnote{This flexibility, achievable via the in-kernel transpose, eliminates the necessity for shuffle instructions to alter thread register ownership, as was previously necessary in~\citep{colfax_fp8_flashattention_2024}.}

\textbf{Accuracy: block quantization and incoherent processing.} The FP8 (e4m3) representation allocates only 3 bits to the mantissa and 4 bits to the exponent, leading to greater numerical inaccuracies compared to FP16 or BF16. Additionally, large-scale models are known to produce outlier values~\citep{dettmers2208llm, sun2024massive}, whose magnitudes far exceed those of most elements, complicating the quantization process. A standard practice is to adopt per-tensor scaling~\citep{micikevicius2022fp8}, in which a separate scaling factor is maintained for each tensor (for example, individual scalars for $\mathbf{Q}$, $\mathbf{K}$, and $\mathbf{V}$).
To mitigate the numerical imprecision in FP8 attention, we introduce two strategies:

\begin{enumerate}

\item \textbf{Block quantization}: In this approach, we maintain a single scalar for each block. Specifically, for every tensor among $\mathbf{Q}$, $\mathbf{K}$, and $\mathbf{V}$, we partition the tensor into blocks of dimension $B_r \times d$ or $B_c \times d$ and quantize each block independently. This quantization procedure can be seamlessly integrated with operations that immediately precede attention (such as rotary embedding), incurring no additional latency, as such steps are typically constrained by memory bandwidth. Given that the \textsc{FlashAttention-3} algorithm inherently works on block structures, it is straightforward to apply a distinct scaling factor to each block in $\mathbf{S}$ to compensate for block quantization without incurring any extra computational overhead.

\item \textbf{Incoherent processing}: To mitigate the impact of outliers, we pre-multiply both $\mathbf{Q}$ and $\mathbf{K}$ with a random orthogonal matrix $\mathbf{M}$ prior to performing FP8 quantization. Because $\mathbf{M}$ is orthogonal, it satisfies $\mathbf{M} \mathbf{M}^\top = I$, and thus $(\mathbf{Q} \mathbf{M}) (\mathbf{K} \mathbf{M})^\top = \mathbf{Q} \mathbf{K}^\top$, ensuring that applying $\mathbf{M}$ does not alter the resulting attention computation. The application of $\mathbf{M}$ effectively disperses any outlier values, as each component of $\mathbf{Q} \mathbf{M}$ or $\mathbf{K} \mathbf{M}$ becomes a random linear combination of the original entries, thereby diminishing quantization errors. Following the methodology in \citet{chee2024quip} and \citet{tseng2024quip}, we construct $\mathbf{M}$ as the product of random diagonal matrices containing $\pm 1$ entries and a Hadamard matrix. This structure enables matrix multiplication in $O(d \log d)$ time, rather than the usual $O(d^2)$, and this operation can also be combined with rotary embedding at no additional computation cost.
\end{enumerate}

We validate that these two techniques reduces numerical error by up to 2.6$\times$ in \cref{sec:numerical_error}.

\section{Empirical Validation}
\label{sec:exp}

We use the primitives from CUTLASS~\citep{Thakkar_CUTLASS_2023} such as WGMMA and TMA abstractions to implement \textsc{FlashAttention-3} and evaluate its efficiency and accuracy.

\begin{itemize}

\item \textbf{Benchmarking attention.} We evaluate the execution time of \textsc{FlashAttention-3} over a range of sequence lengths and contrast its performance against the baseline PyTorch implementation, \textsc{FlashAttention-2}, a Triton-based version of \textsc{FlashAttention-2} (leveraging H100-specific capabilities), as well as a proprietary variant of \textsc{FlashAttention-2} from cuDNN tailored for H100 GPUs. Our results indicate that \textsc{FlashAttention-3} achieves speedups of up to 2.0$\times$ compared to \textsc{FlashAttention-2} and 1.5$\times$ over the Triton implementation. Additionally, \textsc{FlashAttention-3} attains a peak throughput of 740 TFLOPs/s, corresponding to 75\% of the H100 GPUs’ theoretical maximum.
\item \textbf{Ablation study.} Through targeted ablation experiments, we verify that enhancements such as warp-specialization and GEMM-softmax pipelining significantly drive the performance gains observed in \textsc{FlashAttention-3}.
\item \textbf{Accuracy of FP8 attention.} We demonstrate that techniques like block quantization and incoherent processing decrease the numerical error in FP8 \textsc{FlashAttention-3} computations by a factor of 2.6$\times$.
\end{itemize}

\subsection{Benchmarking Attention}
\label{subsec:benchmark_attn}

We evaluate the execution times of various attention mechanisms on an H100 80GB SXM5 GPU, across multiple configurations (both with and without a causal mask, and using head dimensions of 64 or 128) for FP16 input representations. Our findings, displayed in~\cref{fig:benchmark_attn_fwd} and~\cref{fig:benchmark_attn_bwd}, demonstrate that \textsc{FlashAttention-3} achieves forward pass speeds approximately 1.5-2.0$\times$ greater than those of \textsc{FlashAttention-2}, and backward pass speeds about 1.5-1.75$\times$ faster. When benchmarked against conventional attention implementations, \textsc{FlashAttention-3} exhibits performance gains of up to 3-16$\times$. Notably, for sequence lengths of 1k tokens or more, \textsc{FlashAttention-3} outperforms a highly optimized proprietary vendor library (cuDNN), designed specifically for H100 GPUs.

\paragraph{Benchmark settings:} We adjust the sequence length across 512, 1k, up to 16k, selecting batch sizes such that each batch comprises 16k tokens in total. The hidden dimension is fixed at 2048, with head dimension values of 64, 128, or 256, corresponding to 32, 16, or 8 attention heads, respectively. For forward pass FLOPs computation, we use:

\begin{equation*}
  4 \cdot \text{seqlen}^2 \cdot \text{head dimension} \cdot \text{number of heads}.
\end{equation*}

Applying causal masking reduces the total count by half, as roughly 50\% of the operations are performed. For the backward pass, we estimate its FLOPs by multiplying the forward pass count by 2.5, reflecting that the forward computation involves 2 matrix multiplications, whereas the backward computation—because of recomputation—requires 5.

We additionally benchmark the FP8 forward pass under matching conditions. Results for headdim 256 are shown in ~\cref{fig:benchmark_attn_fp8}, while a comprehensive set of results is presented in \cref{sec:benchmark_attn_fp8_full}.

\subsection{Ablation Study: 2-Stage Pipelining Experiments}
\label{sec:2-stage-pipelining-experiments}

We conduct an ablation study on both the 2-stage WGMMA-softmax pipeline and the application of warp-specialization, focusing on the non-causal FP16 \textsc{FlashAttention-3} with fixed parameters $\{\text{batch}, \text{seqlen}, \text{nheads}, \text{hdim}\} = \{4, 8448, 16, 128\}$. As shown in~\cref{table:ablation_pipelining}, the findings indicate that our enhancements—specifically, introducing asynchrony via warp-specialization and overlapping the GEMM with the softmax computation—produce a notable increase in throughput, raising performance from 570 up to 661 TFLOPs.

\subsection{Numerical Error Validation}
\label{sec:numerical_error}

Given the ongoing discussion about the numerical inaccuracies associated with \textsc{FlashAttention}~\citep{golden2024flash}, we conduct a comparison between \textsc{FlashAttention-2}, \textsc{FlashAttention-3}, and a conventional attention implementation, using a FP64 reference as the baseline. To mimic the presence of outlier features and activations commonly observed in LLMs~\citep{dettmers2208llm, sun2024massive}, we draw the elements of $\mathbf{Q}, \mathbf{K}, \mathbf{V}$ from the distribution described below:

\begin{equation*}
  \mathcal{N}(0, 1) + \mathcal{N}(0, 100) \cdot \mathrm{Bernoulli}(0.001).
\end{equation*}

In this setup, all entries follow a normal distribution centered at zero with a standard deviation of 1, except for 0.1\% of the entries where an independent Gaussian noise term with a standard deviation of 10 is introduced. The root mean squared error (RMSE) is reported in~\cref{table:numerical_error}. When using FP16 precision, \textsc{FlashAttention-2} and \textsc{FlashAttention-3} produce RMSE values that are 1.7$\times$ lower than the conventional implementation, attributed to the retention of intermediate (softmax) values in FP32. For FP8, the baseline attention mechanism employs per-tensor scaling, accumulates the matrix multiplication in FP32, and maintains softmax intermediates in FP16. In comparison, through the adoption of block quantization and incoherent processing, \textsc{FlashAttention-3} in FP8 outperforms the baseline by achieving 2.6$\times$ greater accuracy.

\section{Dicussion, Limitations, Conclusion}
\label{sec:discussion}

Through the development of \textsc{FlashAttention-3}, we illustrate that advances in programming strategies and the adoption of hardware innovations—including asynchrony and lower-precision computation—can substantially enhance both the speed and reliability of attention mechanisms. Our approach achieves an attention speedup of 1.5–2.0$\times$ over \textsc{FlashAttention-2}, while reducing FP8 numerical error by a factor of 2.6 when compared with conventional per-tensor quantization approaches. Several aspects of our current work remain to be improved: further optimizing the framework for LLM inference, integrating a persistent kernel mechanism within the FP8 implementation,\footnote{Our benchmarking reveals that FP16 \textsc{FlashAttention-3} employs a persistent kernel and an advanced load balancing method, whereas FP8 \textsc{FlashAttention-3} does not, contributing in part to the relatively lower performance of FP8 \textsc{FlashAttention-3} on short sequences and with causal masking compared to the FP8 cuDNN kernels.} and more deeply examining the impact of low-precision computation on large-scale training. While this study centers on Hopper GPUs, we believe that the methods described are broadly applicable to a variety of hardware accelerators. We anticipate that increasingly efficient and accurate attention primitives will help enable the development of novel applications for long-context machine learning tasks.

\end{document}